\chapter{Consideraciones Finales}\label{cap:Considerações}

\section{Contribución y alcance}

Esta tesis aporta la primera caracterización que combina densidad espacial (PDFe), EOF univariado y MEOF para cuantificar, por fase del ciclo de vida y región de génesis, el desacople estructural entre viento y precipitación en los ciclones extratropicales del Atlántico Sudoccidental. El hallazgo central --que la intensidad dinámica del vórtice deja de predecir linealmente la magnitud y la ubicación de sus impactos de superficie en el subconjunto p90-- y su respaldo estadístico y geográfico ya se desarrollaron en el Capítulo~\ref{ch:conclusiones}. La confianza en que estos patrones espaciales no son artefactos del muestreo se sustenta en el remuestreo Bootstrap-t (Sección~\ref{subsec:bootstrap_significancia}), que documenta áreas estadísticamente significativas en EOF1 tanto en la Población Global como en el subconjunto p90, extendiéndose en este último también a EOF2 y EOF3 --los mismos modos que capturan la reorganización de mesoescala en los sistemas intensos, en algunos casos con mayor área significativa que el propio EOF1--, lo que respalda que la señal reportada es física y recurrente, y no ruido estocástico. Las limitaciones y líneas futuras que siguen delimitan el alcance de esa contribución.

\section{Limitaciones}

Los resultados están sujetos a la subestimación de ERA5 en extremos ($-10.2\%$ en viento y $-22.6\%$ en precipitación para los ciclones más intensos), por lo que las magnitudes absolutas del subconjunto p90 son estimaciones conservadoras. Las estructuras de mesoescala responsables de los máximos de viento en ciclones severos, como la CCB, los \textit{sting jets} o las descargas descendentes, operan a escalas menores que la resolución de ERA5 ($\sim$31~km), limitando la resolución de su geometría interna en los prototipos. La estandarización del dominio lagrangiano a $\pm 9{,}5^{\circ}$ puede truncar la banda de precipitación enroscada en los ciclones de mayor escala horizontal, subestimando la extensión real del campo pluviométrico periférico. El tamaño muestral del subconjunto p90 difiere sustancialmente entre regiones (198 ciclones en ARG, frente a 64 en SBR y 67 en LPB), por lo que las densidades periféricas de las PDFe y los autovectores del EOF en SBR y LPB están sujetos a mayor incertidumbre de muestreo que los de ARG. La evaluación de significancia estadística mediante Bootstrap-t se concentró en la fase de madurez, el momento de mayor interés para el argumento central de la tesis, de modo que los patrones EOF de las fases restantes se presentan como caracterización descriptiva sin la misma validación formal, y el modo multivariado (MEOF) se apoya en la significancia ya establecida de sus componentes univariados en madurez sin una prueba multivariada independiente. El propio EOF, al ser una descomposición lineal, podría diluir en un solo modo diferencias que en realidad correspondan a subpoblaciones de ciclones estructuralmente discretas, aunque el aplanamiento del espectro de autovalores observado en los sistemas p90 sugiere que el método está capturando, y no ocultando, esa complejidad. El análisis se circunscribe a los campos de superficie. La extensión a la columna tridimensional permitiría vincular las estructuras identificadas con sus mecanismos generadores.

El catálogo de trayectorias empleado \citep{gramcianinov2020analysis, coutodesouza2024new} constituye una única fuente de detección y seguimiento entre varios esquemas objetivos disponibles para el Atlántico Sur (Sección~\ref{sec:tb_vorticidad}); aunque basado en vorticidad relativa a 850~hPa, un criterio ampliamente validado, no puede descartarse que otro catálogo modifique la composición exacta de las poblaciones PG y p90. Se espera, no obstante, que las señales estructurales aquí documentadas --el desacople de cuadrantes, la fragmentación del espectro de autovalores en p90-- sean robustas frente a esta elección, en la medida en que dependen de la física del ciclón y no del algoritmo de detección particular.

A diferencia de estudios del Hemisferio Norte que rotan cada composite para alinear la dirección de desplazamiento del sistema, por ejemplo hacia el este, antes de promediar \citep{rudeva2011composite, priestley2022improved}, el marco lagrangiano aquí empleado mantiene la orientación geográfica original (Sección~\ref{subsec:tb_lagrangiano}). Esta elección preserva la interpretación directa de los cuadrantes en términos cardinales (noroeste, sureste), pero no corrige la variabilidad en la dirección de traslación de los ciclones, lo que podría diluir, nunca exagerar artificialmente, las asimetrías estructurales reportadas; que estas emerjan con claridad pese a esta fuente adicional de dispersión respalda la robustez de la señal detectada.

Adicionalmente, la exigencia de un ciclo de vida completo con las cuatro fases y sin re-intensificaciones (Sección~\ref{subsec:grupo_global}) excluye sistemas de evolución compleja, truncada o de vida corta, que podrían representar una fracción no despreciable de los ETC reales de la región. Los resultados aquí obtenidos caracterizan, por tanto, el subconjunto de ciclones con evolución estructuralmente completa, no la totalidad de la actividad ciclónica del Atlántico Sudoccidental.

\section{Perspectivas futuras}

\begin{enumerate}
    \item \textit{Validación con observaciones independientes.} La geometría de los prototipos debe contrastarse con observaciones de alta resolución (boyas, estaciones costeras, altimetría satelital) para cuantificar los sesgos de ERA5 y ajustar los valores modales de las PDFe en el subconjunto p90, en particular la extensión real de la banda enroscada de precipitación.

    \item \textit{Extensión tridimensional.} Incorporar los campos de viento en 850~hPa, temperatura potencial equivalente y vorticidad potencial al análisis MEOF permitirá vincular las firmas superficiales con la evolución tridimensional de las corrientes transportadoras (WCB, CCB) y la estructura del \textit{sting jet} en columna.

    \item \textit{Análisis de eventos compuestos.} Los prototipos proveen el marco de referencia para definir umbrales combinados sobre las PDFe de viento y precipitación, cuantificando la probabilidad de impactos compuestos en las costas de Brasil, Uruguay y Argentina con diferenciación explícita entre el peligro eólico (noroeste) y el hidrológico (sureste-sur) en función de la posición del ciclón respecto a la costa.

    \item \textit{Proyecciones climáticas.} Evaluar en simulaciones de alta resolución con escenarios SSP si el calentamiento global modifica la separación de cuadrantes documentada, en particular la posición y extensión del núcleo de viento noroeste y la coherencia de la banda enroscada, con implicaciones directas para la gestión del riesgo costero en el Atlántico Sudoccidental.

    \item \textit{Extensión de la validación estadística.} Aplicar el remuestreo Bootstrap-t a las cuatro fases del ciclo de vida y directamente al modo multivariado (MEOF), y ampliar el registro histórico o incorporar salidas de modelos de alta resolución para SBR y LPB, permitiría reducir la incertidumbre de muestreo respecto a ARG y fortalecer la robustez estadística de los prototipos estructurales en las regiones con menor tamaño muestral.
\end{enumerate}
